% FORMALIZACION_REUNIDA: representación pentacomponente de II,
% con antecedente Paley y demostración directa de los momentos del marco.
% El retorno transporta la carta; no se identifica un ciclo de orden doce con A5.
\section{Incidencia orientada y representación tensorial gravitatoria}
\label{vii:sec:pentacomponente}
La representación excepcional del estado conserva la incidencia que acompaña
a su lectura aritmética. La matriz ternaria \(A_W\) de
\eqref{exc:aw} proporciona aquí un antecedente material preciso. Su
levantamiento balanceado \(C=\operatorname{bal}(A_W)\), con representantes
\(0,1,-1\), satisface
\[
 C=C^{\mathsf T},\qquad C^2=5I_6,\qquad \operatorname{Tr}C=0.
\]
Las seis posiciones tienen el orden \((\infty,0,1,2,3,4)\).
Esta incidencia produce primero un espacio tridimensional; su lectura
cuadrática produce después el espacio de cinco componentes sobre el que
actúa el acoplamiento gravitatorio.

\subsection{Recubrimiento firmado y transporte del calendario}
Definimos
\[
 E_C=\frac12(I_6+C/\sqrt5),\quad
 \mathcal H_C=\operatorname{im}E_C,\quad
 v_{u,\sigma}=\sigma\sqrt2E_Ce_u,\quad \sigma\in\{1,-1\}.
\]
\begin{proposition}[Doce posiciones orientadas]
\label{vii:prop:doce-orientadas}
Los vectores \(v_{u,\sigma}\) son doce vectores unitarios distintos
en un espacio de dimensión tres. Su incidencia está determinada por
\(\langle v_{u,\sigma},v_{v,\tau}\rangle=1/\sqrt5\).
La inversión de la orientación actúa por antipodía.
\end{proposition}
\begin{proof}
La identidad cuadrática de \(C\) implica \(E_C^2=E_C=E_C^*\);
su traza es tres. Puesto que la diagonal de \(C\) es nula,
\[
 \langle v_{u,\sigma},v_{v,\tau}\rangle
 =\begin{cases}\sigma\tau,&u=v,\\
 \sigma\tau C_{uv}/\sqrt5,&u\ne v.
 \end{cases}
\]
Esto prueba normalización, distinción de los doce vectores e incidencia.
Cambiar \(\sigma\) por \(-\sigma\) cambia el signo del vector.
\end{proof}

Una expresión euclidiana explícita se obtiene poniendo
\(\nu=\sqrt{1+\varphi^2}\) y
\[
 B=\frac1\nu
 \begin{pmatrix}
 0&-1&-\varphi&0&\varphi&1\\
 -1&-\varphi&0&1&0&-\varphi\\
 -\varphi&0&-1&-\varphi&-1&0
 \end{pmatrix}.
\]
La relación \(\varphi^2=\varphi+1\), para la coordenada ya generada,
da \(B^*B=2E_C\).
En consecuencia \(B/\sqrt2|_{\mathcal H_C}\) es una isometría y lleva
\(v_{u,\sigma}\) a \(\sigma Be_u\). La imagen es
\[
 \mathcal V=\frac1\nu\bigl(
 \{0\}\times\{\pm1\}\times\{\pm\varphi\}\ \cup\
 \{\pm1\}\times\{\pm\varphi\}\times\{0\}\ \cup\
 \{\pm\varphi\}\times\{0\}\times\{\pm1\}\bigr).
\]
Así quedan relacionados la coordenada interna, la incidencia y sus
doce representantes geométricos.

La coordenada del calendario se escribe \((b,r)\in
\mathbb Z/12\mathbb Z\times\{0,\ldots,8\}\). Una ruta de longitud \(n\)
induce
\[
 k(r,n)=\left\lfloor\frac{r+n}{9}\right\rfloor,\qquad
 (b,r)\longmapsto
 \bigl(b+k(r,n)\bmod12,\ (r+n)\bmod9\bigr).
\]
El alineamiento explícito de sus posiciones \(p=b+1\) con la coordenatización firmada es
\[
\begin{array}{c|rrrrrrrrrrrr}
p&1&2&3&4&5&6&7&8&9&10&11&12\\ \hline
u&\infty&0&0&\infty&1&3&3&1&2&4&4&2\\
\sigma&+&-&+&-&-&+&-&+&-&+&-&+
\end{array}
\]
Esta tabla fija una trivialización del portador. Cada cambio de posición
transporta conjuntamente los vectores y su matriz de incidencia.

\begin{proposition}[Covariancia del marco bajo retorno]
Si \(S e_p=e_{p+1\bmod12}\), la representación matricial de un
operador \(P\) del portador se transporta como \(P_k=S^kPS^{-k}\).
Este transporte conserva productos escalares, incidencias y composición
de rutas. Tras doce retornos vuelve la coordenatización, conservándose la memoria
acumulada en el estado fuente.
\end{proposition}
\begin{proof}
\(S\) es ortogonal y \(S^{12}=I\); de aquí se siguen las propiedades
métricas y la periodicidad de la coordenatización. Para la composición, con
\(r'=(r+n)\bmod9\), la división euclidiana da
\(k(r,n+m)=k(r,n)+k(r',m)\). Los exponentes de \(S\) se suman
exactamente con el acarreo. La operación actúa sobre la coordenada
del calendario; las actualizaciones de memoria de la ruta siguen
presentes en su dominio enriquecido.
\end{proof}

\subsection{Lectura cuadrática y espacio de cinco componentes}
Sea \(\mathbb R^{\mathcal V}\) el espacio de las doce coordenadas,
\(\mathsf A e_v=e_{-v}\) su involución antipodal y
\(\mathsf J=\mathbf1\mathbf1^{\mathsf T}\). Se definen
\[
 P_5=\frac12(I+\mathsf A)-\frac1{12}\mathsf J,\qquad
 Q_v=vv^{\mathsf T}-\frac13I_3,\qquad
 \Gamma_0x=\sum_{v\in\mathcal V}x_vQ_v.
\]
\begin{theorem}[Proyector pentacomponente e isometría tensorial]
\label{vii:thm:portador-pentacomponente}
\(P_5\) es el proyector ortogonal de rango cinco sobre el espacio de
coordenadas pares bajo antipodía y de suma nula. Además,
\[
 \Gamma_0^*\Gamma_0=\frac85P_5,\qquad
 \Gamma_0\Gamma_0^*=\frac85I_{\operatorname{Sym}^2_0(\mathbb R^3)}.
\]
Por tanto
\[
 \Gamma_5=\sqrt{\frac58}\,\Gamma_0|_{\operatorname{im}P_5}
\]
es una isometría sobre el espacio de tensores simétricos sin traza.
\end{theorem}
\begin{proof}
El espacio par tiene una coordenada por cada par antipodal, por lo que
su dimensión es seis. El vector constante pertenece a él; restar su
proyector de rango uno da \(P_5\), de rango cinco.

Como
\(\langle Q_v,Q_w\rangle_F=(v^{\mathsf T}w)^2-1/3\), las entradas
de \(\Gamma_0^*\Gamma_0\) valen \(2/3\) para \(w=\pm v\) y
\(-2/15\) en los otros diez lugares de cada fila. Éstas son
exactamente las entradas de \((8/5)P_5\).

También puede comprobarse la identidad en el espacio tensorial sin
recurrir a una clasificación de representaciones. La lista de vértices da
\[
 \sum_vv_i^4=\frac{12}{5},\qquad
 \sum_vv_i^2v_j^2=\frac45\quad(i\ne j),
\]
y los momentos con alguna potencia impar se anulan por los cambios de
signo. En efecto,
\((1+\varphi^2)^2=5\varphi^2\) y \(1+\varphi^4=3\varphi^2\).
Para un tensor simétrico \(Q\), estas sumas dan
\[
 \sum_v(v^{\mathsf T}Qv)\,vv^{\mathsf T}
 =\frac85Q+\frac45(\operatorname{tr}Q)I.
\]
Con \(\operatorname{tr}Q=0\), y usando
\(\sum_vv^{\mathsf T}Qv=4\operatorname{tr}Q=0\), resulta
\(\Gamma_0\Gamma_0^*Q=(8/5)Q\). Las dos identidades prueban
la isometría normalizada y su sobreyectividad.
\end{proof}

La sección gravitatoria de \(\ref{v:ant:sec:gravity}\), con la
normalización demostrada en~\ref{g26:sec:rigidez-radial-es}, determina
el operador \(G_aP_5\). En la sección de retorno,
\(G_a=c^3L^2/\hbar_{\rm ret}\). La isometría tensorial transporta
este mismo operador mediante
\[
 \Gamma_5(G_aP_5)|_{\operatorname{im}P_5}\Gamma_5^*
   =G_a I_{\operatorname{Sym}^2_0(\mathbb R^3)}.
\]
La igualdad se deduce de \(P_5|_{\operatorname{im}P_5}=I\) y
\(\Gamma_5\Gamma_5^*=I\). Así, el factor \(\sqrt{5/8}\) fija la
normalización tensorial y la composición radial fija el coeficiente
escalar, con sus respectivas demostraciones.
El espectro del operador \(G_aP_5\) sobre \(\mathbb R^{12}\)
tiene \(G_a\) con multiplicidad
cinco y cero con multiplicidad siete. La orientación se conserva en
el portador y en su transporte; \(G_a\) fija el acoplamiento común.

\subsection{Proyección transversal y helicidades}
Para \(n\in S^2\), pongamos \(\Pi=I-nn^{\mathsf T}\) y
\[
 \Lambda_nQ=\Pi Q\Pi-\frac12\operatorname{tr}(\Pi Q\Pi)\Pi.
\]
\begin{proposition}[Reducción transversal de rango dos]
\(\Lambda_n\) es el proyector ortogonal sobre
\(\{Q\in\operatorname{Sym}^2_0(\mathbb R^3):Qn=0\}\), de dimensión dos.
\end{proposition}
\begin{proof}
\(\Pi n=0\) da transversalidad y \(\operatorname{tr}\Pi=2\) da
traza nula. Un tensor transversal sin traza queda fijo.
La fórmula es autoadjunta para el producto de Frobenius.
En un marco \((e_1,e_2,n)\), la imagen tiene base ortonormal
\[
 E_+=\frac{e_1e_1^{\mathsf T}-e_2e_2^{\mathsf T}}{\sqrt2},\qquad
 E_\times=\frac{e_1e_2^{\mathsf T}+e_2e_1^{\mathsf T}}{\sqrt2}.
\]
\end{proof}
Los otros tres vectores de una base ortonormal del portador son
\[
 E_1=\frac{e_1n^{\mathsf T}+ne_1^{\mathsf T}}{\sqrt2},\quad
 E_2=\frac{e_2n^{\mathsf T}+ne_2^{\mathsf T}}{\sqrt2},\quad
 E_0=\frac{2nn^{\mathsf T}-e_1e_1^{\mathsf T}-e_2e_2^{\mathsf T}}{\sqrt6}.
\]
Al rotar el marco transversal un ángulo \(\theta\), las parejas
\((E_+,E_\times)\) y \((E_1,E_2)\) giran \(2\theta\) y
\(\theta\), respectivamente, mientras \(E_0\) queda fijo. La
complejificación tiene así los pesos \(2,-2,1,-1,0\).
\(\Lambda_n\) conserva los dos primeros y anula los restantes.
Queda establecida la cadena
\[
 \mathbb R^{12}\xrightarrow{P_5}\operatorname{im}P_5
 \xrightarrow{\Gamma_5}\operatorname{Sym}^2_0(\mathbb R^3)
 \xrightarrow{\Lambda_n}\mathcal H_{\rm TT}(n),
 \qquad 12\longrightarrow5\longrightarrow5\longrightarrow2.
\]
La representación pentacomponente y su reducción son resultados
algebraicos anteriores a la elección de ecuaciones de campo.
En la realización gravitatoria sin masa, las restricciones de calibre
seleccionan la imagen transversal. Esta distinción conserva tanto las
cinco coordenadas estructurales como las dos helicidades radiativas.
